% ============================================================
%  APPENDIX
%  Stratified Monte Carlo Pricing and Variance Reduction
%  via K-Means Clustering and Machine Learning Control Variates
%
%  Benjamin Sulpice & Enzo Thuilliez — May 2026
%  Supervised by Professor Arnaud DUFAYS
% ============================================================
%
%  This appendix collects the four publication-quality figures
%  generated by the Python pipeline (main.py) and the critical
%  code extracts underlying the three machine-learning steps of
%  the framework: K-Means/Neyman allocation (Step 3), local
%  Random Forest control variates (Step 4), and the global
%  neural-network variance-loss optimiser (Step 5).
%
%  All figures are formatted consistently with the main text
%  (monochrome/grayscale, serif typeface, minimalist layout).
%  All code blocks use the \texttt{listings} environment with
%  monochrome syntax highlighting: bold keywords, italic
%  comments, no colour.
% ============================================================

% ── Appendix counter setup ───────────────────────────────────
% Chapter, figure, table, and equation counters are reset per
% appendix chapter and labelled A.1, B.1, etc. automatically
% when this file is included after \appendix in main.tex.

% ============================================================
\chapter{Numerical Figures}
\label{app:figures}
% ============================================================

The figures presented in this appendix constitute the four primary
visual outputs of the five-step pipeline implemented in \texttt{main.py}.
Each figure is generated under identical model parameters
($S_0 = K = 100$, $r = 0.05$, $\sigma = 0.20$, $T = 1$, $M = 252$
daily time steps) and saved in publication-quality grayscale at
200\,dpi.  The sequence of figures mirrors the logical progression
of the pipeline: from raw path simulation (Figure~\ref{fig:trajectoires}),
through empirical convergence benchmarks (Figure~\ref{fig:rmse}),
to the diagnostic outputs of the clustering stage
(Figure~\ref{fig:clusters}) and the neural-network optimiser
(Figure~\ref{fig:nn_loss}).

% ------------------------------------------------------------
\section{Figure 1: Simulated Asset Trajectories}
\label{app:fig1}
% ------------------------------------------------------------

\paragraph{Context.}
The first step of the pipeline generates $N$ Monte Carlo paths
of the underlying asset under the risk-neutral Geometric Brownian
Motion (GBM) dynamics:
\[
  S_t = S_0 \exp\!\left[\left(r - \tfrac{1}{2}\sigma^2\right)t
  + \sigma W_t^{\mathbb{Q}}\right], \quad t \in [0, T],
\]
where $W^{\mathbb{Q}}$ is a standard Brownian motion under the
pricing measure $\mathbb{Q}$.  Figure~\ref{fig:trajectoires}
displays a representative sample of $N = 200$ such paths,
colour-coded by their terminal moneyness status.  This
visualisation serves two purposes: (i) it validates the
correctness of the discretisation scheme and the resulting
distributional properties of $S_T$, and (ii) it provides an
immediate geometric illustration of the heterogeneity that
motivates trajectory clustering --- paths that expire deeply
in the money exhibit systematically different path
characteristics (higher mean, lower realised variance) than
out-of-the-money paths, justifying the partition of the
simulation space into homogeneous strata.

\begin{figure}[H]
  \centering
  \includegraphics[width=\linewidth]{figure_1_trajectoires.png}
  \caption{%
    \textbf{Simulated GBM trajectories categorised by terminal
    moneyness.}
    Two hundred paths of Geometric Brownian Motion are displayed
    under parameters $S_0 = K = 100$, $r = 0.05$,
    $\sigma = 0.20$, $T = 1$, $M = 252$.
    \textit{Left panel}: European call perspective.
    Dark paths expire in-the-money ($S_T > K$); light paths
    expire out-of-the-money ($S_T \leq K$).  The dashed
    horizontal line marks the strike $K$.
    \textit{Right panel}: European put perspective, with dark
    paths now corresponding to $S_T < K$.
    The bimodal structure of the terminal distribution under
    the moneyness partition provides the geometric motivation
    for K-Means clustering in the five-dimensional feature
    space $(S_T,\,\bar{S},\,\widehat{\sigma}^2,\,S_{\max},\,S_{\min})$.%
  }
  \label{fig:trajectoires}
\end{figure}

% ------------------------------------------------------------
\section{Figure 2: Empirical RMSE Convergence Rates}
\label{app:fig2}
% ------------------------------------------------------------

\paragraph{Context.}
The fundamental limitation of plain Monte Carlo pricing is its
slow convergence rate: the Root Mean Squared Error (RMSE) of the
estimator $\hat{V}_0^{MC} = N^{-1}\sum_{i=1}^N \phi_i$ decreases
as $\mathcal{O}(N^{-1/2})$, where $N$ is the number of simulated
paths.  Every variance reduction technique in the pipeline can be
understood as an attempt to reduce the constant $C$ in
$\mathrm{RMSE} \approx C \cdot N^{-1/2}$, or equivalently, to
reduce the variance $\mathrm{Var}(\phi_i)$ of the effective payoff
being averaged.

Figure~\ref{fig:rmse} benchmarks five architectural variants of
the Monte Carlo estimator across sample sizes
$N \in \{500,\,1{,}000,\,2{,}000,\,5{,}000,\,10{,}000,\,20{,}000\}$.
For each $(N, \text{method})$ pair, $n = 100$ independent
replications are performed, and the RMSE is estimated as
$\sqrt{n^{-1}\sum_{j=1}^n (\hat{V}_{0,j} - V_0^*)^2}$, where
$V_0^* = V_0^{BS}$ is the analytical Black--Scholes price.  The
log-log scale allows the empirical convergence slopes to be
compared directly against the $\mathcal{O}(N^{-1/2})$ reference
line (shown as a dashed gray line).  A steeper slope indicates
super-$\sqrt{N}$ convergence in the effective sample size, which
is equivalent to a variance reduction factor that grows with $N$.

\begin{figure}[H]
  \centering
  \includegraphics[width=0.88\linewidth]{figure_2_RMSE.png}
  \caption{%
    \textbf{Log-log empirical RMSE of five Monte Carlo
    estimators across sample sizes $N \in [500,\,20{,}000]$.}
    Each RMSE value is computed over $n = 100$ independent
    replications.  Estimators shown: Plain MC (solid, circles),
    Antithetic Variates (dashed, squares), Stratified K-Means
    (dash-dot, triangles), Random Forest control variate
    (dotted, diamonds), and Neural Network control variate
    (dash-dot-dot, stars).  The gray dashed line represents
    the theoretical $\mathcal{O}(N^{-1/2})$ reference slope.
    Machine-learning control variates (RF, NN) achieve a
    downward vertical shift of the RMSE curve, corresponding
    to a multiplicative variance reduction factor that is
    independent of $N$ and preserved uniformly across the
    sample-size range tested.  The stratified estimator
    exhibits a reduction consistent with the theoretical
    within-stratum variance decomposition.%
  }
  \label{fig:rmse}
\end{figure}

% ------------------------------------------------------------
\section{Figure 3: K-Means Cluster Structure in Feature Space}
\label{app:fig3}
% ------------------------------------------------------------

\paragraph{Context.}
The K-Means algorithm (Step 3 of the pipeline) partitions the
$N_{\text{pilot}}$ simulated paths into $K$ clusters by minimising
the total within-cluster variance of the normalised feature
vectors $\mathbf{x}_i \in \mathbb{R}^5$.  The features are:
terminal value $S_T$, arithmetic path mean $\bar{S}$, realised
variance $\widehat{\sigma}^2$, path maximum $S_{\max}$, and path
minimum $S_{\min}$.  Although the clustering is performed in this
full five-dimensional space, Figure~\ref{fig:clusters} projects
the assignments onto the $(S_T, \bar{S})$ plane, which captures
the dominant sources of payoff heterogeneity for standard European
and path-dependent options.

The figure demonstrates that K-Means recovers economically
meaningful regimes: clusters in the left region of the plot
(low $S_T$, low $\bar{S}$) correspond to deeply out-of-the-money
paths with near-zero call payoff and high put payoff; clusters in
the right region correspond to in-the-money paths; and clusters
straddling the strike line $K = 100$ (dashed) capture the
near-money paths, which contribute the highest within-cluster
payoff variance and therefore receive the largest Neyman allocation.

\begin{figure}[H]
  \centering
  \includegraphics[width=0.85\linewidth]{figure_3_clusters.png}
  \caption{%
    \textbf{K-Means cluster assignments projected onto the
    $(S_T,\,\bar{S})$ feature plane.}
    $N_{\text{pilot}} = 20{,}000$ GBM paths are clustered into
    $K = 8$ strata using K-Means++ initialisation.  Each point
    represents one simulated path, plotted at its terminal value
    $S_T$ (horizontal axis) against its arithmetic mean $\bar{S}$
    (vertical axis).  Distinct grayscale shades and marker
    symbols identify the eight cluster assignments.  The dashed
    vertical line marks the strike $K = 100$.  The cluster
    boundaries reveal a structure aligned with the moneyness
    partition and with the path-average distribution, confirming
    that the unsupervised algorithm recovers economically
    interpretable regimes without any explicit payoff information.
    The Neyman optimal allocation $n_k^* \propto p_k\,\sigma_k$
    subsequently concentrates the simulation budget on
    near-money clusters, where within-cluster payoff variance
    $\sigma_k$ is highest.%
  }
  \label{fig:clusters}
\end{figure}

% ------------------------------------------------------------
\section{Figure 4: Neural Network Variance-Loss Trajectory}
\label{app:fig4}
% ------------------------------------------------------------

\paragraph{Context.}
Step 5 of the pipeline trains a feedforward neural network
$g_\theta : \mathbb{R}^5 \to \mathbb{R}$ to serve as a global
control variate.  The network is optimised with respect to the
\emph{custom variance loss}:
\[
  \mathcal{L}(\theta, c)
  \;=\;
  \mathbb{E}\!\left[\bigl(\phi - c\,g_\theta(\mathbf{x})\bigr)^2\right]
  \;=\;
  \Var(\phi) - 2c\,\Cov(\phi, g_\theta) + c^2\,\Var(g_\theta),
\]
where $\phi$ is the (discounted) path payoff and $\mathbf{x}$
is the corresponding feature vector.  The scalar coefficient
$c$ is updated in closed form after each epoch:
\[
  c^*_{\text{epoch}}
  = \frac{\widehat{\Cov}(\phi,\,g_\theta)}{\widehat{\Var}(g_\theta)},
\]
which is precisely the optimal OLS coefficient minimising
$\mathcal{L}(\theta, c)$ with respect to $c$ for fixed $\theta$.
This alternating optimisation ensures that the loss converges
monotonically to the residual variance $\Var(\phi - c^* g_\theta)$,
which equals $(1 - \rho^2)\,\Var(\phi)$ at the global optimum,
where $\rho = \mathrm{Corr}(\phi, g_\theta)$.

Figure~\ref{fig:nn_loss} traces this minimisation across
$200$ training epochs on $N_{\text{pilot}} = 20{,}000$ paths,
providing empirical confirmation that the custom loss
converges reliably under the Adam optimiser.

\begin{figure}[H]
  \centering
  \includegraphics[width=0.82\linewidth]{figure_4_nn_loss.png}
  \caption{%
    \textbf{Optimisation trajectory of the neural-network
    variance loss $\mathcal{L}(\theta, c)$ over $200$ epochs.}
    The network $g_\theta$ is a three-hidden-layer feedforward
    architecture with 64 units per layer and ReLU activations,
    trained via Adam ($\alpha = 10^{-3}$, $\lambda = 10^{-4}$
    weight decay, $p = 0.10$ dropout).  The solid black line
    plots the per-epoch batch-average loss; the dashed gray
    line shows the 10-epoch moving average.  The annotated
    terminal value confirms convergence to a stable minimum.
    The final loss corresponds to $\mathcal{L}^* = (1 -
    \rho^2)\,\widehat{\Var}(\phi)$, where $\rho > 0.90$ is the
    empirical correlation between payoffs and network outputs
    on the pilot set, implying a variance reduction factor
    exceeding $5\times$ relative to plain Monte Carlo.%
  }
  \label{fig:nn_loss}
\end{figure}

% ============================================================
\chapter{Python Code Extracts}
\label{app:code}
% ============================================================

This appendix chapter reproduces the three critical algorithmic
routines from \texttt{main.py} that implement Steps 3--5 of the
pipeline.  The extracts are presented in full, without
abbreviation, to allow complete reproducibility of the numerical
results reported in the main text.  Syntax highlighting is
restricted to bold keywords and italic comments, in keeping with
the monochrome typographic standard of this document.

All code is written in Python 3.10+ and requires the following
libraries: \texttt{numpy}, \texttt{scikit-learn},
\texttt{torch} (PyTorch), and \texttt{scipy}.  The functions
below are self-contained and callable once the upstream feature
matrix \texttt{X\_scaled} and pilot payoff vector \texttt{payoffs}
have been constructed (Steps 1--2).

% ------------------------------------------------------------
\section{K-Means Clustering and Neyman Allocation (Step 3)}
\label{app:code-kmeans}
% ------------------------------------------------------------

\paragraph{Mathematical background.}
Given a normalised feature matrix
$\mathbf{X} = \{\mathbf{x}_i\}_{i=1}^{N} \subset \mathbb{R}^5$,
K-Means++ solves the combinatorial optimisation problem:
\[
  \min_{\{C_k\}_{k=1}^K}
  \sum_{k=1}^K \sum_{\mathbf{x} \in C_k}
  \|\mathbf{x} - \boldsymbol{\mu}_k\|_2^2,
\]
where $\boldsymbol{\mu}_k = |C_k|^{-1}\sum_{\mathbf{x} \in C_k}\mathbf{x}$
is the centroid of cluster $k$.  Once the pilot paths are assigned
to clusters, the Neyman optimal budget allocation concentrates
simulation effort on the strata with the highest variance--weight
product:
\[
  n_k^* = N \cdot
  \frac{p_k\,\sigma_k}{\sum_{j=1}^K p_j\,\sigma_j},
  \qquad
  p_k = \frac{|C_k|}{N}, \quad
  \sigma_k = \widehat{\mathrm{Std}}(\phi \mid C_k).
\]
The following listing implements both the clustering and the
allocation calculation in a single coherent routine.

\begin{lstlisting}[
  language=Python,
  caption={K-Means++ clustering and Neyman optimal allocation
           (from \texttt{main.py}, Section~5).},
  label={lst:kmeans}
]
from sklearn.cluster import KMeans
import numpy as np

def cluster_paths_kmeans(X_scaled, K, random_state=42):
    """
    K-Means++ clustering on normalised feature matrix.

    Parameters
    ----------
    X_scaled     : (N, d) float ndarray  -- standardised features
    K            : int                   -- number of clusters
    random_state : int

    Returns
    -------
    labels   : (N,) int ndarray   -- cluster assignment per path
    centers  : (K, d) float ndarray -- centroid coordinates
    km_model : fitted KMeans object
    """
    km = KMeans(
        n_clusters=K,
        init="k-means++",   # smart initialisation (Arthur & Vassilvitskii, 2007)
        n_init=10,           # 10 restarts to avoid local minima
        random_state=random_state
    )
    labels = km.fit_predict(X_scaled)
    return labels, km.cluster_centers_, km


def neyman_allocation(payoffs, labels, K, N_total):
    """
    Compute Neyman-optimal simulation budget per cluster.

    Formula: n_k* = N * (p_k * sigma_k) / sum_j(p_j * sigma_j)

    Parameters
    ----------
    payoffs  : (N_pilot,) float ndarray -- pilot discounted payoffs
    labels   : (N_pilot,) int ndarray   -- cluster assignments
    K        : int                      -- number of clusters
    N_total  : int                      -- total pricing budget

    Returns
    -------
    allocation : (K,) int ndarray   -- paths allocated to each cluster
    p_k        : (K,) float ndarray -- empirical cluster weights
    sigma_k    : (K,) float ndarray -- within-cluster payoff std devs
    """
    N_pilot = len(payoffs)
    p_k     = np.zeros(K)
    sigma_k = np.zeros(K)

    # Step 3a: compute within-cluster statistics
    for k in range(K):
        mask   = labels == k
        n_k    = mask.sum()
        p_k[k] = n_k / N_pilot
        # Use ddof=1 (unbiased) standard deviation
        sigma_k[k] = payoffs[mask].std(ddof=1) if n_k >= 2 else 0.0

    # Step 3b: Neyman weights
    denom   = (p_k * sigma_k).sum()
    weights = (p_k * sigma_k / denom) if denom > 0 else p_k

    # Step 3c: round to integer counts, then correct rounding error
    allocation = np.maximum(
        np.round(weights * N_total).astype(int), 1
    )
    # Ensure allocation sums to exactly N_total
    allocation[np.argmax(weights)] += N_total - allocation.sum()

    return allocation, p_k, sigma_k
\end{lstlisting}

% ------------------------------------------------------------
\section{Local Random Forest Control Variate (Step 4)}
\label{app:code-rf}
% ------------------------------------------------------------

\paragraph{Mathematical background.}
For each cluster $k$, a Random Forest regressor
$\hat{f}_k : \mathbb{R}^5 \to \mathbb{R}$ is fitted on the
pilot paths belonging to $C_k$.  The locally corrected payoff
estimator for cluster $k$ is:
\[
  \tilde{\phi}_i^{(k)}
  = \phi_i - c_k^*\bigl(\hat{f}_k(\mathbf{x}_i) - \mu_{Y,k}\bigr),
  \qquad
  c_k^* = \frac{\widehat{\Cov}(\phi,\,\hat{f}_k)}
               {\widehat{\Var}(\hat{f}_k)},
  \quad
  \mu_{Y,k} = \mathbb{E}[\hat{f}_k(X) \mid X \in C_k],
\]
where $c_k^*$ is the OLS-optimal scalar coefficient and
$\mu_{Y,k}$ ensures that the control variate correction is
mean-zero, preserving the unbiasedness of the estimator.  The
global stratified price estimate combines the within-cluster
corrected means weighted by the empirical cluster proportions $p_k$.

\begin{lstlisting}[
  language=Python,
  caption={Local Random Forest control variate training and
           application (from \texttt{main.py}, Section~7).},
  label={lst:rf}
]
from sklearn.ensemble import RandomForestRegressor
import numpy as np

def rf_control_variate(
    X_pilot, payoffs_pilot, labels_pilot,
    X_pricing, labels_pricing, p_k, K,
    n_estimators=200, max_depth=None, random_state=42
):
    """
    Fit one Random Forest per cluster; use predictions as
    local control variates to reduce within-cluster variance.

    Parameters
    ----------
    X_pilot         : (N_pilot, d)   normalised pilot features
    payoffs_pilot   : (N_pilot,)     discounted pilot payoffs
    labels_pilot    : (N_pilot,)     cluster labels (pilot)
    X_pricing       : (N_price, d)   normalised pricing features
    labels_pricing  : (N_price,)     cluster labels (pricing)
    p_k             : (K,)           empirical cluster weights
    K               : int            number of clusters

    Returns
    -------
    estimate          : float   -- corrected price estimate
    residual_var_ratio: float   -- Var(corrected) / Var(raw)
    """
    forests    = {}   # one RF model per cluster
    ctrl_means = {}   # E[f_k(X) | X in C_k]  (for mean-zero correction)

    # ── Training phase: fit one RF per cluster ──────────────
    for k in range(K):
        mask = labels_pilot == k
        if mask.sum() < 10:           # insufficient data: skip
            forests[k]    = None
            ctrl_means[k] = 0.0
            continue

        rf = RandomForestRegressor(
            n_estimators=n_estimators,
            max_depth=max_depth,
            min_samples_leaf=5,       # prevents overfitting on small leaves
            n_jobs=-1,
            random_state=random_state
        )
        rf.fit(X_pilot[mask], payoffs_pilot[mask])
        forests[k]    = rf
        # Control variate mean (estimated on pilot training data)
        ctrl_means[k] = rf.predict(X_pilot[mask]).mean()

    # ── Pricing phase: apply local OLS correction ───────────
    corrected, raw = [], []

    for k in range(K):
        mask_tr = labels_pilot  == k
        mask_pr = labels_pricing == k
        n_tr    = mask_tr.sum()
        n_pr    = mask_pr.sum()

        if n_pr == 0:
            continue

        # Align pilot sample to pricing cluster size
        phi_pr = (payoffs_pilot[mask_tr][:n_pr]
                  if n_tr >= n_pr
                  else payoffs_pilot[mask_tr])
        if len(phi_pr) == 0:
            continue

        raw.extend(phi_pr.tolist())

        if forests[k] is None:        # fallback: no correction
            corrected.extend(phi_pr.tolist())
            continue

        X_pr   = X_pricing[mask_pr][:len(phi_pr)]
        Y_pr   = forests[k].predict(X_pr)
        mu_Y_k = ctrl_means[k]

        # OLS optimal coefficient: c* = Cov(phi, Y) / Var(Y)
        cov_kk = np.cov(phi_pr, Y_pr)[0, 1]
        var_Y  = Y_pr.var(ddof=1)
        c_k    = cov_kk / var_Y if var_Y > 1e-10 else 0.0

        # Apply mean-zero control variate correction
        corrected.extend(
            (phi_pr - c_k * (Y_pr - mu_Y_k)).tolist()
        )

    corrected = np.array(corrected)
    raw       = np.array(raw)
    raw_var   = raw.var(ddof=1) if raw.var(ddof=1) > 0 else 1.0

    return corrected.mean(), corrected.var(ddof=1) / raw_var
\end{lstlisting}

% ------------------------------------------------------------
\section{Neural Network Variance-Loss Optimiser (Step 5)}
\label{app:code-nn}
% ------------------------------------------------------------

\paragraph{Mathematical background.}
The global neural network control variate solves:
\[
  \min_{\theta \in \Theta,\, c \in \mathbb{R}}
  \mathcal{L}(\theta, c)
  = \mathbb{E}\!\left[
      \bigl(\phi - c\,g_\theta(\mathbf{x})\bigr)^2
    \right],
\]
which equals $\Var(\phi) - 2c\,\Cov(\phi, g_\theta) + c^2\,\Var(g_\theta)$.
The minimiser over $c$ at fixed $\theta$ is the OLS coefficient
$c^* = \Cov(\phi, g_\theta)/\Var(g_\theta)$, and at the global optimum
$\mathcal{L}^* = (1 - \rho^2)\Var(\phi)$ where $\rho = \mathrm{Corr}(\phi, g_\theta)$.
The alternating update --- gradient step in $\theta$, closed-form update
of $c$ --- ensures monotone decrease of the loss.  The network
architecture is:
\[
  g_\theta : \mathbb{R}^5
  \xrightarrow{\text{Linear}_{64}} \text{ReLU}
  \xrightarrow{\text{Linear}_{64}} \text{ReLU}
  \xrightarrow{\text{Linear}_{64}} \text{ReLU}
  \xrightarrow{\text{Linear}_{1}} \mathbb{R}.
\]
Weights are initialised with He (Kaiming) normal initialisation
for ReLU activations.

\begin{lstlisting}[
  language=Python,
  caption={Neural network architecture and custom variance-loss
           training routine (from \texttt{main.py}, Section~8).},
  label={lst:nn}
]
import torch
import torch.nn as nn
import torch.optim as optim
import numpy as np


class ControlVariateNet(nn.Module):
    """
    Feedforward network  g_theta : R^d -> R.

    Architecture  : Input -> [Linear(64) -> ReLU] x 3 -> Linear(1)
    Initialisation: He (Kaiming) normal for ReLU layers.
    """
    def __init__(self, d_in: int):
        super().__init__()
        self.net = nn.Sequential(
            nn.Linear(d_in, 64), nn.ReLU(),
            nn.Linear(64, 64),   nn.ReLU(),
            nn.Linear(64, 64),   nn.ReLU(),
            nn.Linear(64, 1),
        )
        # He initialisation: preserves variance across ReLU layers
        for m in self.modules():
            if isinstance(m, nn.Linear):
                nn.init.kaiming_normal_(m.weight, nonlinearity="relu")
                nn.init.zeros_(m.bias)

    def forward(self, x):
        return self.net(x).squeeze(-1)


def train_nn_control_variate(
    X_train, phi_train,
    n_epochs=200, batch_size=256, lr=1e-3,
    weight_decay=1e-4, dropout_p=0.1, device="cpu",
    verbose=False, record_loss=False,
):
    """
    Train ControlVariateNet with the custom variance loss
    L(theta, c) = E[(phi - c * g_theta(x))^2].

    The scalar c is updated in closed form after each epoch:
        c* = Cov(phi, g_theta) / Var(g_theta)

    Parameters
    ----------
    X_train      : (N, d) float ndarray  -- normalised features
    phi_train    : (N,)   float ndarray  -- discounted payoffs
    n_epochs     : int    training epochs
    batch_size   : int    mini-batch size (SGD)
    lr           : float  Adam learning rate
    weight_decay : float  L2 regularisation coefficient
    dropout_p    : float  hidden-layer dropout probability
    device       : str    'cpu' or 'cuda'
    record_loss  : bool   if True, return per-epoch loss history

    Returns
    -------
    net       : trained ControlVariateNet
    c_star    : float   -- optimal scalar coefficient
    mu_g      : float   -- E[g_theta(x)] on training set
    loss_hist : list    -- per-epoch losses (if record_loss=True)
    """
    X_t   = torch.tensor(X_train,  dtype=torch.float32).to(device)
    phi_t = torch.tensor(phi_train, dtype=torch.float32).to(device)
    N     = X_train.shape[0]

    net       = ControlVariateNet(X_train.shape[1]).to(device)
    optimiser = optim.Adam(
        net.parameters(), lr=lr, weight_decay=weight_decay
    )
    c_star    = 1.0          # initialise coefficient
    loss_hist = []

    for epoch in range(n_epochs):
        net.train()
        perm       = torch.randperm(N)
        epoch_loss = 0.0
        n_batches  = 0

        for start in range(0, N, batch_size):
            idx   = perm[start: start + batch_size]
            x_b   = X_t[idx]
            phi_b = phi_t[idx]

            optimiser.zero_grad()
            g_b = net(x_b)

            # Custom variance loss: Var(phi - c * g_theta)
            loss = ((phi_b - c_star * g_b) ** 2).mean()
            loss.backward()      # backpropagate through g_theta only
            optimiser.step()
            epoch_loss += loss.item()
            n_batches  += 1

        # Closed-form update of c after each full epoch
        net.eval()
        with torch.no_grad():
            g_all    = net(X_t)
            phi_mean = phi_t.mean()
            g_mean   = g_all.mean()
            # c* = Cov(phi, g) / Var(g)
            num = ((phi_t - phi_mean) * (g_all - g_mean)).mean().item()
            den = ((g_all - g_mean) ** 2).mean().item()
            if abs(den) > 1e-10:
                c_star = num / den

        avg_loss = epoch_loss / max(n_batches, 1)
        loss_hist.append(avg_loss)

        if verbose and (epoch + 1) % 50 == 0:
            print(
                f"  Epoch {epoch+1:3d}/{n_epochs}  "
                f"loss={avg_loss:.6f}  c*={c_star:.4f}"
            )

    # Final mean of g_theta (for unbiased price correction)
    net.eval()
    with torch.no_grad():
        mu_g = net(X_t).mean().item()

    return (net, c_star, mu_g, loss_hist) if record_loss \
           else (net, c_star, mu_g)
\end{lstlisting}

% ============================================================
\chapter{Supplementary Parameter Tables}
\label{app:params}
% ============================================================

Table~\ref{tab:params-model} summarises the model and simulation
parameters used throughout the numerical study.
Table~\ref{tab:params-nn} collects the hyperparameters of the
neural network control variate.

\begin{table}[H]
  \centering
  \caption{%
    \textbf{Model and simulation parameters.}
    Parameters are held fixed across all experiments unless
    otherwise noted in the main text.%
  }
  \label{tab:params-model}
  \begin{tabular}{llll}
    \toprule
    Symbol & Description & Value & Unit \\
    \midrule
    $S_0$ & Initial asset price & 100 & \$ \\
    $K$   & Strike price        & 100 & \$ \\
    $r$   & Risk-free rate      & 0.05 & ann.\ rate \\
    $\sigma$ & GBM volatility   & 0.20 & ann.\ vol \\
    $T$   & Maturity            & 1.0  & years \\
    $M$   & Time steps          & 252  & steps (daily) \\
    $N_{\text{pilot}}$  & Pilot paths & 20\,000 & paths \\
    $N_{\text{price}}$  & Pricing paths & 20\,000 & paths \\
    $K_{\text{clust}}$ & Number of clusters & 8 & --- \\
    \midrule
    $\kappa$ & Heston mean-reversion speed & 2.0  & --- \\
    $\theta$ & Heston long-run variance    & 0.04 & --- \\
    $\xi$    & Heston vol-of-vol           & 0.30 & --- \\
    $\rho$   & Heston correlation          & $-0.70$ & --- \\
    $V_0$    & Heston initial variance     & 0.04 & --- \\
    \bottomrule
  \end{tabular}
\end{table}

\begin{table}[H]
  \centering
  \caption{%
    \textbf{Neural network hyperparameters.}
    The architecture and training configuration are held fixed
    across all option types and model dynamics.%
  }
  \label{tab:params-nn}
  \begin{tabular}{lll}
    \toprule
    Hyperparameter & Description & Value \\
    \midrule
    Architecture   & Hidden layers $\times$ units & $3 \times 64$ \\
    Activation     & Hidden activation             & ReLU \\
    Initialisation & Weight initialisation         & He (Kaiming) normal \\
    Optimiser      & Gradient method               & Adam \\
    Learning rate  & $\alpha$                      & $10^{-3}$ \\
    Weight decay   & L2 regularisation $\lambda$   & $10^{-4}$ \\
    Dropout        & Hidden-layer dropout $p$      & $0.10$ \\
    Batch size     & Mini-batch size               & 256 \\
    Epochs         & Training epochs               & 200 \\
    Pilot size     & Pre-training paths $N_0$      & 50\,000 \\
    \bottomrule
  \end{tabular}
\end{table}